\chapter{Current limitations}
\label{chap:limitations}

\section{Thin vessels are below the resolution}
\label{sec:thin}

With a blur of $\sigma = 4.4$\um{}, the image of a point is about 10\um{} wide at half its height. A
vessel of 3 to 6\um{} radius is narrower than the blur: its image is essentially the blur's own shape,
dimmer than a large vessel, whatever its true size. On the phantom, vessels of that size are found
\pbValueVesselRecallThin{} of the time, against \pbValueVesselRecallMedium{} for 6--12\um{} and
\pbValueVesselRecallLarge{} above.

I tried two recoveries and dropped them. The missed thin vessels turned out not to be faint at all (their
voxels are mostly brighter than the bone threshold): they are lost because, next to a brighter vessel,
they fall below that vessel's half-maximum and are taken for its halo. Telling them apart from a real
halo with the Hessian recovered about as many wrong voxels as right ones: no net gain, so I did not
keep the code. The
consequence for future measurements is that diameters under about 12\um{} will be overestimated, and
the thinnest vessels partly missed: this is a property of the scan, not of the software. The optional
thin-vessel pass (Section~\ref{sec:thinvessels}) undoes part of the blur and finds
\pbValueExpThinThinRecall{} of them on the phantom instead of \pbValueVesselRecallThin; what it adds
on the real scan is unconfirmed (Section~\ref{sec:experimental}).

\section{Rods and vessels}

Shape alone cannot tell a trabecular rod from a vessel: both are tubes. The decision rests on what they
touch in 3D. A dim vessel running along a trabecula, or a rod touching vessels on all sides, can still
be misclassified; these are the same ``zones of uncertainty'' we met in 2018.

\section{Validation}

The phantom is calibrated on the real scan, but it is synthetic: its vessel paths, burst textures and
bone geometry are simplifications. The scores of Chapter~\ref{chap:validation} measure the method on
realistic data, not the segmentation of this particular knee. Expert labels on a few real crops (the
annotation page is ready) would close that gap.

\section{The data}

\begin{itemize}
  \item \textbf{One sample.} I could only check the parameters on a single scan. Physical units make them
        transferable in principle, not in fact.
  \item \textbf{Voxel size.} 3\um{} is the slice thickness we used in 2018 and agrees with the
        anatomy; I assumed isotropy.
  \item \textbf{A thin slab.} The stack is 100 slices thick (0.3\,mm) for a knee several millimetres
        wide: most vessels leave the volume, which limits what 3D connectivity can decide near the
        first and last slices.
  \item \textbf{Bursts and vessels.} A vessel running into a burst cluster is partly absorbed by it,
        and the edge of a burst takes the bright core of a thick vessel passing close by.
\end{itemize}

\section{Other scans}

The thresholds are computed from each scan and every size is in micrometres, so another resolution or
brightness range should only need configuration changes. Other limits are built into the method:
\begin{itemize}
  \item \textbf{Three intensity classes.} The method assumes background $<$ bone $<$ contrast agent:
        the vessels must be the brightest material (barium or another strongly radio-opaque agent) and
        the bone must not be decalcified. A preparation where the agent looks dimmer than bone would
        need new code, not new settings.
  \item \textbf{One set of thresholds per stack.} A scan whose brightness drifts along the stack is
        classified with thresholds that fit neither end.
  \item \textbf{Fixed rules.} A vessel is a bright core with a rim, a tube rather than a plate; bone is
        large and connected; a burst is an area dense in bright speckles. The numbers behind these rules
        are settings; the rules are not.
  \item \textbf{Values tuned on one sample.} Minimum sizes, burst density and the tube-versus-plate
        cut-off were chosen on this scan and on a phantom calibrated on it. Without a reference, nothing
        tells whether they still fit another scan.
  \item \textbf{Format.} One 16-bit grey-level image per slice, ordered by file name.
\end{itemize}
A public scan would rarely come with vessel annotations: segmenting it would give one more sample to
look at, not a measured accuracy.

\section{The web server}

The server has no login and must stay on the local machine; it reads stacks from its own disk (no
upload), and keeps its jobs in memory: stopping it cancels the running ones and forgets the others.
Meshes go to the browser as GLB (Table~\ref{tab:meshfiles}): \pbValueMeshVesselGlbMB\,MB for the
vessels of the sample, \pbValueMeshBoneGlbMB\,MB for the bone. A recent graphics card draws both, but
they fill its memory, and the \pbValueMeshGridUm\um{} grid is already a compromise: full resolution
would quadruple it. Nothing is compressed: on
the local machine, compressing would take longer than sending. The page needs a recent browser
(JavaScript modules, WebGL~2). A job asked for while the same one runs is joined, whatever its
settings: other settings apply to the next one.

\section{The vascular network}

The brief asked for a quantified network: nodes, branches, calibre, branch lengths. pyBones now
measures it (Chapter~\ref{chap:network}), within the limits measured there: lengths are the net of
three errors (the measure reads long, thin vessels are missed, strands follow the holes of the vessel
mask); junction counts are inflated by the loops the skeleton draws where the vessel mask has holes or
folds (most junctions of the segmented phantom are not forks), mostly where a vessel touches bone. The
setting that measures a vessel's edge against bone on the bone side brings them from
\pbValueExpDefaultSpurious{} to \pbValueExpBoneSideSpurious{} on the branching phantom, but stays off
until annotated crops of the real scan can confirm it (Section~\ref{sec:experimental}). Where two
vessels pass closer than the blur's width (about 10\um), the image holds one bright shape and the
network joins them: cutting along the dimmest path between their cores (a watershed), with or
without a test of the image's curvature, removed only a few of these junctions, at the cost of real
forks and length, and was dropped. The radius is biased in opposite directions for thin and
thick vessels; and many
branches are truncated by the thin stack. The measures cover the whole stack: no region (cortex,
marrow) is told apart yet.

\section{Two experimental switches}
\label{sec:experimental}

Two settings improve the phantoms but stay off by default, because what they change on the real scan
cannot be checked without expert labels: measuring a vessel's edge against bone on the bone side of a
contact (\code{segmentation.vessel.bone\_side}, Section~\ref{sec:edges}), and adding the thin vessels a
deconvolved image shows (\code{thin\_vessels}, Section~\ref{sec:thinvessels}).
Table~\ref{tab:experimental} measures them on both phantoms, alone and together; the report tool runs
them every time.

\begin{table}[ht]
\centering
\small
\begin{tabular}{lrrrrrrr}
\toprule
 & \multicolumn{3}{c}{Standard phantom (vessels)} & \multicolumn{4}{c}{Branching phantom} \\
\cmidrule(lr){2-4}\cmidrule(lr){5-8}
Setup & Dice & Precision & Recall $<6$\um & Forks & Other & Loops & Length \\
\midrule
Both off (default) & 0.862 & 0.826 & 0.43 & 26/53 & 49 & 7 & $-3.0$\,\% \\
Bone-side edges & 0.887 & 0.876 & 0.42 & 30/53 & 28 & 5 & $-7.2$\,\% \\
Thin vessels & 0.863 & 0.826 & 0.60 & 27/53 & 50 & 7 & $+2.1$\,\% \\
Both on & 0.888 & 0.875 & 0.60 & 31/53 & 30 & 5 & $-1.8$\,\% \\
\bottomrule
\end{tabular}

\caption{The experimental switches on the standard phantom (vessel scores, recall of vessels thinner
  than 6\um) and on the branching phantom (forks found of those visible, other junctions, independent
  loops, error of the total length).}
\label{tab:experimental}
\end{table}

Their effects add up: the bone-side edges remove the false vessel along bone, which removes many of
the junctions that are not forks, but shorten the network; the thin vessels give most of that length
back. On the real scan, however, the bone-side edges change the vessel volume by
\pbValueExpRealBoneSideVesselPct\,\%, turning into bone thin, dim structures next to bone, and the
thin vessels change it by \pbValueExpRealThinVesselPct\,\%, as bright spots and short bars in or
along bone. Both are plausible, and neither can be confirmed: a few annotated crops of the real scan,
compared side by side, would settle them. Until then they can be switched on per run, from the
command line or the web page's settings form, where they come first under \emph{Experimental}.
